% !TeX program = pdflatex
%
% pk-research-paper.tex
%
% Build:  pdflatex pk-research-paper.tex   (run 3x: ToC, then cross-refs)
%

\documentclass[11pt,a4paper]{article}

\usepackage[T1]{fontenc}
\usepackage[utf8]{inputenc}
\usepackage[margin=2.4cm]{geometry}
\usepackage{amsmath}
\usepackage{amssymb}
\usepackage{booktabs}
\usepackage{longtable}
\usepackage{array}
\usepackage{xcolor}
\usepackage{enumitem}
\usepackage{fancyhdr}
\usepackage[hidelinks]{hyperref}

\definecolor{warn}{RGB}{160,60,0}
\definecolor{ok}{RGB}{20,110,60}
\newcommand{\bad}[1]{\textcolor{warn}{#1}}
\newcommand{\good}[1]{\textcolor{ok}{#1}}
\newcommand{\code}[1]{\texttt{\small #1}}

\pagestyle{fancy}
\fancyhf{}
\lhead{\small xT Pharmacokinetic Model}
\rhead{\small\thepage}
\renewcommand{\headrulewidth}{0.4pt}

\title{\textbf{Superposed Bateman Biexponentials with Clearance-Constrained\\
Scaling and Per-User Calibration:}\\[0.3em]
\large A Pharmacokinetic Framework for Hormone Replacement\\
Therapy Monitoring in the xT Mobile Application}
\author{Starla Cortes}
\date{August 2026}

\begin{document}
\maketitle

\begin{abstract}
\noindent
We present the pharmacokinetic (PK) model underlying xT, a mobile
application that predicts serum hormone concentration over time from a
user's logged dose history for gender-affirming hormone replacement
therapy (HRT). The model employs a one-compartment Bateman
biexponential---first-order absorption from an intramuscular depot into a
single well-mixed compartment, first-order elimination from it---with
linear superposition of all prior dose contributions, a dynamic
endogenous baseline that decays under gonadal suppression, and a
per-user calibration factor $\alpha$ derived from self-reported blood
tests. Ten medication/route/form combinations are supported: estradiol
(valerate, cypionate, oral, transdermal) and testosterone (cypionate,
enanthate, propionate, undecanoate, oral, transdermal).

The scaling constant $S$---an empirical surrogate for $F/V_d$---is
derived in closed form for testosterone cypionate from a primary
population-PK source (Bi et al.\ 2018), and for the remaining injectable
testosterone esters by a conservation-of-mass identity that constrains
$S$ to depend only on the elimination rate $k_e$ and testosterone
clearance, independent of the ester's half-life. Estradiol ester
constants are anchored to compiled steady-state PK tables. Transdermal
routes are modeled as 24 hourly micro-doses approximating square-wave
release.

We report accuracy against published anchors (0.0--0.1\% for anchored
quantities), out-of-sample validation (1--15\% for estradiol and
testosterone single-dose and steady-state checks), and time-to-peak
($t_{\max}$) agreement with primary and secondary sources. We
quantify structural limitations: testosterone undecanoate's peak/trough
ratio (8.1$\times$ modeled vs.\ 1.8$\times$ observed, a source conflict
rather than a structural deficiency), transdermal intraday ripple
($C_{\max}/C_{\min} \approx 1.05$ vs.\ 2.24 observed), and oral
testosterone peak fluctuation under-prediction. The model is
honest about where its constants are primary-sourced, where they are
secondary, and where they remain uncalibrated.

\textbf{Keywords:} pharmacokinetics, Bateman function, hormone
replacement therapy, testosterone cypionate, estradiol valerate,
flip-flop kinetics, superposition, per-user calibration, mobile health
\end{abstract}

\tableofcontents
\clearpage

% ======================================================================
\section{Introduction}
\label{sec:intro}

Gender-affirming hormone replacement therapy (HRT) involves the
administration of exogenous hormones---estradiol and/or
testosterone---to align an individual's serum hormone levels with their
gender identity. Clinical management relies on periodic blood draws to
measure serum hormone concentration, titrate dosage, and confirm
therapeutic levels~\cite{hembree2017}. Between draws, patients and
providers have limited tools to anticipate where levels are heading,
when a peak or trough is expected, or how a missed dose or regimen
change will affect the trajectory.

xT is a mobile application that addresses this gap by predicting serum
hormone concentration over time from the user's logged dose history. It
implements a pharmacokinetic model that computes, for any point in time,
the superposed contribution of every dose the user has taken, plus a
dynamic endogenous baseline, scaled by an optional per-user calibration
factor derived from blood test results.

The model covers ten medication/route/form combinations across two
hormone classes (Table~\ref{tab:scope}). Its mathematical core is the
Bateman biexponential---the analytic solution to one-compartment
first-order absorption and elimination---which the published
pharmacometric literature independently identifies as the best-fit model
for every injectable testosterone ester examined: cypionate~\cite{bi2018},
enanthate~\cite{ichihara2020}, and
undecanoate~\cite{pastuszak2021}. This convergence is not incidental;
it reflects the fact that depot testosterone esters, across a wide range
of chain lengths and dosing intervals, produce terminal phases that are
single-exponential on semi-log plots, justifying one-compartment fits.

This paper describes the model in full: its governing equation, the
provenance and derivation of every constant, the calibration mechanism,
the accuracy of its predictions against published data, and the
structural limitations that bound what it can honestly claim. It is
written to be checkable: every constant is stated with its source, and
every derivation is given in enough detail to reproduce.

\subsection{Contributions}

\begin{enumerate}[nosep,leftmargin=1.4em]
\item A complete description of a pharmacokinetic model designed for
  consumer HRT tracking, including the full constant set, per-form
  derivations, and the dynamic endogenous baseline model.
\item A conservation-of-mass identity (Section~\ref{sec:identity}) that
  constrains the scaling constant $S$ for all injectable testosterone
  esters to depend only on $k_e$ and clearance, eliminating $S$ as a
  free parameter and enforcing mutual consistency across esters.
\item An analytic closed-form derivation of $S$ for testosterone
  cypionate from primary population-PK parameters
  (Section~\ref{sec:tc}), which serves as the anchor for the identity
  above.
\item A per-user calibration framework ($\alpha = C_{\text{measured}} /
  C_{\text{predicted}}$) that closes part of the between-subject
  variability gap without requiring a formal mixed-effects model.
\item An honest accounting of structural limitations, including a
  source-conflict diagnosis for testosterone undecanoate's peak/trough
  ratio that reframes a previously assumed structural deficiency as a
  calibration problem.
\end{enumerate}

% ======================================================================
\section{Methods}
\label{sec:methods}

\subsection{System overview}
\label{sec:overview}

xT predicts serum hormone concentration over time from a user's logged
dose history. Two medication classes are supported, each across three
administration routes, yielding ten medication/route/form combinations
(Table~\ref{tab:scope}).

\begin{table}[h]
\centering
\small
\begin{tabular}{@{}llll@{}}
\toprule
\textbf{Medication} & \textbf{Route} & \textbf{Ester / form} & \textbf{Output unit} \\
\midrule
Estradiol     & Injectable   & Valerate (EV), Cypionate (EC) & pg/mL \\
Estradiol     & Oral         & (route pseudo-ester)          & pg/mL \\
Estradiol     & Cream        & (route pseudo-ester)          & pg/mL \\
Testosterone  & Injectable   & Cypionate (TC), Enanthate (TE), & ng/dL \\
              &              & Propionate (TP), Undecanoate (TU) & \\
Testosterone  & Oral         & undecanoate (route pseudo-ester) & ng/dL \\
Testosterone  & Cream        & unesterified gel (route pseudo-ester) & ng/dL \\
\bottomrule
\end{tabular}
\caption{The ten medication/route/form combinations the model covers.}
\label{tab:scope}
\end{table}

\paragraph{Route pseudo-esters.} The two non-injectable cases in each
ester enumeration describe a \emph{route}, not a chemical ester. They
are resolved from the dose record's administration method by a
route-to-ester mapping function, never chosen by the user directly. This
design ensures that a dose logged as ``oral testosterone'' uses the
oral route's bioavailability and rate constants, not those of whatever
injectable ester was previously selected.

\paragraph{Exogenous and endogenous separation.} The PK engine returns
the \emph{exogenous} contribution only---the sum of all administered
dose contributions. Endogenous production is modeled separately as a
dynamic baseline $C_{\text{base}}(t)$ that decays as gonadal suppression
progresses (Section~\ref{sec:baseline}). This separation matters for
interpreting several constants, most importantly testosterone
cypionate's half-life: the value encoded (3.84\,d) is the
exogenous-only disposition rate, not the total-testosterone terminal
half-life of $\sim$7\,d that is commonly quoted and that includes
endogenous secretion.

\subsection{Governing equation}
\label{sec:governing}

Each administered dose contributes a Bateman biexponential---first-order
absorption from a depot into a single well-mixed compartment, first-order
elimination from it:

\begin{equation}
C_i(t) = D_i \cdot \mathrm{EDF} \cdot S \cdot
         \frac{k_a}{k_a - k_e}\left(e^{-k_e t} - e^{-k_a t}\right),
         \qquad t \ge 0
\label{eq:bateman}
\end{equation}

where $D_i$ is the administered dose in mg, $\mathrm{EDF}$ is the
effective dose factor (Section~\ref{sec:edf}), $S$ is a scaling constant
with units of concentration per mg, and $k_a, k_e$ are rate constants in
day$^{-1}$. Total predicted level at time $t$ is the superposition of
all doses taken before it, plus the endogenous baseline, times a user
calibration factor:

\begin{equation}
C(t) = \alpha\left[C_{\text{base}}(t) + \sum_{i \,:\, t_i \le t} C_i(t - t_i)\right]
\label{eq:superposition}
\end{equation}

Superposition is exact here because Equation~\ref{eq:bateman} is linear
in dose. That linearity is also the model's principal simplification:
there is no saturable protein binding, no SHBG dynamics, no
dose-dependent clearance.

\paragraph{The $k_a \approx k_e$ limit.} When $|k_a - k_e| < 10^{-4}$,
Equation~\ref{eq:bateman} is singular, and the engine substitutes its
limit,
\[
D \cdot \mathrm{EDF} \cdot S \cdot k_a\, t\, e^{-k_a t},
\]
which is the standard degenerate-case form of the Bateman function when
the two rate constants coincide.

\paragraph{Transdermal route.} The cream/gel route does not use
Equation~\ref{eq:bateman}. It is modeled as 24 hourly micro-doses of
$D/24$ each, decaying single-exponentially ($k_e = 0$), approximating a
24-hour square-wave release:
\begin{equation}
C_{\text{transdermal}}(t) = \sum_{h=0}^{23} \frac{D}{24} \cdot S \cdot e^{-k_a (t - t_h)},
\qquad t \ge t_h
\label{eq:transdermal}
\end{equation}
where $t_h$ is the time of the $h$-th micro-dose. This construction has
no absorption phase; Section~\ref{sec:limitations} explains why it
under-models intraday variation.

\subsection{Time to peak}
\label{sec:tmax}

Differentiating Equation~\ref{eq:bateman} with respect to $t$ and
setting the derivative to zero gives the time to peak concentration:

\begin{equation}
t_{\max} = \frac{\ln(k_a/k_e)}{k_a - k_e}
\label{eq:tmax}
\end{equation}

A property of Equation~\ref{eq:tmax} that has practical consequences is
its \emph{symmetry} under $k_a \leftrightarrow k_e$: swapping the two
rate constants leaves $t_{\max}$ unchanged. For depot esters the two
exponentials are not separately identifiable from clinical data, so
which one is labelled ``absorption'' and which ``elimination'' is a
modelling convention rather than a finding. The observable curve is
unchanged by the swap, with the amplitude difference absorbed into $S$.

\subsection{Flip-flop kinetics}
\label{sec:flipflop}

For depot esters, absorption from the oil vehicle is slower than
elimination of the released hormone ($k_a \ll k_e$). The terminal slope
of the curve is then governed by \emph{absorption}, not elimination---
``flip-flop'' kinetics~\cite{rowland2011}. The practical consequence is
that the parameter controlling the terminal half-life sets $k_a$ and
governs the entire inter-dose peak-to-trough decay, while the parameter
controlling the fast exponential sets $k_e$ and primarily determines
where the peak falls.

In xT's parameterization, the field \code{halfLifeDays} sets $k_a = \ln
2 / t_{1/2,\text{term}}$ (the slow rate), and
\code{eliminationHalfLifeHours} sets $k_e = \ln 2 / (t_{1/2,\text{elim}}
/ 24)$ (the fast rate). The naming is historical; the doc comments in
the codebase explicitly note the convention.

\subsection{Effective dose factor}
\label{sec:edf}

$\mathrm{EDF} = \mathrm{EWF} \times F$ converts administered mg of ester
into systemically available mg of active hormone.

\subsubsection{Ester weight factor}

The ester weight factor is pure stoichiometry---the mass fraction of the
ester molecule that is the parent hormone:
\[
\mathrm{EWF} = \frac{M_{\text{hormone}}}{M_{\text{ester}}}
\]
with $M_{\mathrm{E}_2} = 272.4$ and $M_{\mathrm{T}} = 288.4$\,g/mol.
These are the only constants in the entire model that are not fitted to
anything and cannot be wrong. Table~\ref{tab:ewf} gives the values for
all esters.

\begin{table}[h]
\centering
\small
\begin{tabular}{@{}llrl@{}}
\toprule
\textbf{Ester} & \textbf{Abbrev.} & \textbf{EWF} & \textbf{Molecular basis} \\
\midrule
Estradiol valerate    & EV       & 0.764 & $272.4 / 356.5$ \\
Estradiol cypionate   & EC       & 0.687 & $272.4 / 396.6$ \\
Estradiol (oral/TD)   & E2       & 1.000 & base hormone, no ester \\
Testosterone cypionate & TC      & 0.699 & $288.4 / 412.6$ \\
Testosterone enanthate & TE      & 0.720 & $288.4 / 400.6$ \\
Testosterone propionate & TP     & 0.837 & $288.4 / 344.5$ \\
Testosterone undecanoate & TU    & 0.631 & $288.4 / 456.7$ \\
Testosterone (gel)    & T        & 1.000 & base hormone, no ester \\
\bottomrule
\end{tabular}
\caption{Ester weight factors from molecular weights. Oral estradiol
and transdermal routes deliver the base hormone directly, so EWF = 1.}
\label{tab:ewf}
\end{table}

\subsubsection{Bioavailability}

$F$ is route bioavailability---the fraction of the administered dose
reaching systemic circulation:

\begin{table}[h]
\centering
\small
\begin{tabular}{@{}lll@{}}
\toprule
\textbf{Route} & \textbf{$F$} & \textbf{Basis} \\
\midrule
Injectable depot & 1.0 & Near-complete absorption from oil depot \\
Oral estradiol   & 0.05 & First-pass metabolism \\
Transdermal E2   & 0.15 & Skin barrier limits absorption \\
Transdermal T    & 0.10 & AndroGel 1\% label \S12.3 (stated verbatim) \\
Oral TU          & 1.0  & Not established; folded into $S$ (see text) \\
\bottomrule
\end{tabular}
\caption{Route bioavailability values.}
\label{tab:bioavailability}
\end{table}

Oral testosterone undecanoate carries $F = 1.0$ not because it is fully
bioavailable---it is not---but because its absolute bioavailability is
not established in the source. Since $F$ and $S$ are confounded (only
their product is identified by the published data), folding
bioavailability into $S$ avoids inventing a number. Predictions are
unaffected by where the split is drawn.

\subsection{The scaling constant}
\label{sec:scaling}

$S$ is an empirical surrogate for $F/V_d$. It is the one constant with
no direct physical interpretation, and it is the constant most sensitive
to calibration errors. Wherever the sources allow, $S$ is derived in
closed form rather than by curve fit; the derivations are given in
Sections~\ref{sec:tc} and \ref{sec:identity}.

\subsection{Testosterone cypionate: primary-source derivation}
\label{sec:tc}

Testosterone cypionate (TC) is the only form in the model derived from a
primary population-PK publication: Bi et al.~\cite{bi2018}, who report
$\mathrm{CL}/F = 2.6$\,kL/day, $V/F = 14.4$\,kL, and $K_a =
1.22$\,day$^{-1}$ (Table 2 of the paper, with reported relative standard
errors and bootstrap confidence intervals).

\subsubsection{Rate constants}

The disposition rate is $\mathrm{CL}/V = 2.6/14.4 = 0.1806$\,day$^{-1}$,
giving a half-life of 3.839\,d, which sets $k_a$ (the slow rate under
flip-flop). $K_a = 1.22$ gives $\ln 2 / 1.22 = 13.64$\,h, which sets
$k_e$ (the fast rate). Per the symmetry noted after
Equation~\ref{eq:tmax}, the labels are conventional; the observable
curve is unchanged by the swap, with the amplitude difference absorbed
into $S$.

\subsubsection{Why 3.84 days and not 7--8}

The commonly quoted TC half-life is around a week. Bi et al.\ address
this discrepancy directly: noncompartmental analysis of their own data
yields 6.9\,d, and their model yields 6.87\,d \emph{when endogenous
secretion is assumed to be zero}---but 4.05\,d once endogenous production
is accounted for. Because xT models endogenous production separately as
$C_{\text{base}}(t)$, the exogenous-only figure is the one that belongs
in the ester constant. Bi's own median post-hoc estimate of 4.05\,d
agrees with the encoded 3.839\,d to 5.2\%.

\subsubsection{Scaling constant in closed form}

Matching Equation~\ref{eq:bateman} to the one-compartment model
underlying Bi's parameters gives $S$ analytically rather than by curve
fit:
\begin{equation}
S = \frac{10^5 \cdot K_a}{V_L \cdot \mathrm{EDF} \cdot k_e}
  = \frac{10^5 \cdot 1.22}{14400 \times 0.699 \times 0.180556} = 67.13
\label{eq:tcS}
\end{equation}
The $10^5$ converts mg/L to ng/dL; $\mathrm{EDF}$ appears because the
app multiplies the dose by it while Bi's $V/F$ already absorbs ester
conversion. This reproduces Bi's population model exactly.

\subsubsection{Time to peak}

Bi's parameters give $t_{\max} = 1.84$\,d via Equation~\ref{eq:tmax}.
Propagating the published 95\% CI on $K_a$ (0.785--2.67) gives a band of
1.08--2.43\,d. Bi's own limitations section (1--2 absorption-phase
samples per subject, $K_a$ RSE 44.9\%) explains the imprecision. Read TC
$t_{\max}$ as $1.84 \pm 0.7$\,d.

\subsection{Clearance identity for TE, TP, and TU}
\label{sec:identity}

Testosterone enanthate (TE), propionate (TP), and undecanoate (TU) lack
primary population-PK sources with reported $V/F$ and $K_a$ for the
intramuscular route in the model's target population. Their scaling
constants are nonetheless \emph{not} free parameters, because of a
conservation argument that follows from the TC derivation.

\subsubsection{Derivation}

Integrating Equation~\ref{eq:bateman} over all time gives exactly
\[
\int_0^{\infty} \frac{k_a}{k_a-k_e}\left(e^{-k_e t}-e^{-k_a t}\right)\,dt
= \frac{k_a}{k_a-k_e}\left(\frac{1}{k_e}-\frac{1}{k_a}\right) = \frac{1}{k_e}
\]
so at steady state with dosing interval $\tau$, superposition gives
\begin{equation}
C_{\mathrm{avg}} = \frac{D \cdot \mathrm{EDF} \cdot S}{\tau \, k_e}
\label{eq:cavg}
\end{equation}
The true steady-state average is (mg of testosterone per day) divided by
clearance---and clearance is a property of \emph{testosterone}, not of
the ester carrying it. It is therefore already pinned by the TC refit:
$\mathrm{CL}_T = 2600 \times \mathrm{EWF}_{TC} = 1817.4$\,L/day.
Equating the two expressions, the ester weight factor cancels:
\begin{equation}
\boxed{\;S = \frac{10^5 \cdot k_e}{\mathrm{CL}_T \cdot F}\;}
\label{eq:identity}
\end{equation}

$S$ depends only on $k_e$ and clearance. The ester's half-life ($k_a$)
does not appear at all---a property worth stating plainly, because it
means $S$ does not track the ester's terminal half-life but rather its
fast rate constant, which for depot esters under flip-flop is the
absorption rate from the oil vehicle.

\subsubsection{The load-bearing check}

TC satisfies Equation~\ref{eq:identity} to 0.03\% \emph{without having
been fit to it}---its $S$ came from Equation~\ref{eq:tcS}, via $V/F$ and
$K_a$. Two independent derivations landing on the same number is what
makes it defensible to apply the identity to the other three esters.

\subsubsection{Application to TE, TP, and TU}

Table~\ref{tab:identity} shows the identity applied. TE and TU now carry
real absorption rates from primary sources (Section~\ref{sec:te_tu}),
while TP retains a placeholder 2\,h $k_e$ and is flagged as uncalibrated.

\begin{table}[h]
\centering
\small
\begin{tabular}{@{}lrrrr@{}}
\toprule
\textbf{Ester} & $k_e$ (d$^{-1}$) & $S$ (encoded) & $S$ (implied) & \textbf{Error} \\
\midrule
Cypionate   & 1.2196 & 67.13  & 67.11  & \good{0.03\%} \\
Enanthate   & 1.9625 & 107.98 & 107.98 & \good{0.00\%} \\
Propionate  & 8.3178 & 457.67 & 457.67 & \good{0.00\%} \\
Undecanoate & 0.3748 & 20.62  & 20.62  & \good{0.00\%} \\
\bottomrule
\end{tabular}
\caption{The clearance identity (Equation~\ref{eq:identity}) applied to
all injectable testosterone esters. $\mathrm{CL}_T = 1817.4$\,L/day,
$F = 1.0$ for all.}
\label{tab:identity}
\end{table}

\subsubsection{Steady-state predictions at typical regimens}

\begin{table}[h]
\centering
\small
\begin{tabular}{@{}llrrr@{}}
\toprule
\textbf{Ester} & \textbf{Regimen} & $C_{\mathrm{avg}}$ & $C_{\max}$ & $C_{\min}$ \\
\midrule
TC & 100\,mg q7d   & 550 & 736 & 321 \\
TE & 100\,mg q7d   & 566 & 703 & 405 \\
TP & 50\,mg q2d    & 1151 & 1743 & 760 \\
TU & 1000\,mg q70d & 496 & 1045 & 129 \\
\bottomrule
\end{tabular}
\caption{Steady-state predictions (exogenous only, ng/dL) at typical
regimens. $C_{\max}$ and $C_{\min}$ are the peak and trough of the
steady-state oscillation.}
\label{tab:steady_state}
\end{table}

\subsection{Testosterone enanthate and undecanoate: absorption rates}
\label{sec:te_tu}

\subsubsection{Testosterone enanthate (TE)}

TE's absorption rate is derived from Ichihara et al.~\cite{ichihara2020},
the only identified primary IM pharmacokinetic study in transgender men:
$n=10$, single 250\,mg IM dose (Enarmon Depot). The paper reports
per-subject absorption rate constants ($k_a$) in the 4 of 10 subjects
where it was estimable: 2.82, 3.16, 0.49, 1.38\,day$^{-1}$, giving a
mean of 1.9625\,day$^{-1}$ ($t_{1/2} = 8.48$\,h). The terminal half-life
is reported as $6.6 \pm 1.9$\,d, which sets $k_a$ (the slow rate) at
$\ln 2 / 6.6 = 0.1050$\,day$^{-1}$.

The resulting $t_{\max} = 1.58$\,d agrees with the paper's observed
$1.7 \pm 0.9$\,d (all subjects peaked within 3 days). The $t_{\max}$
and the $k_a$ values were estimated in different subsets (Tmax in all
10, $k_a$ in the 4 where it was estimable), so this is a genuine
cross-check rather than a round trip.

\paragraph{Magnitude discrepancy.} A single 250\,mg dose gives 881\,ng/dL
exogenous here against Ichihara's observed $C_{\max}$ of $2940 \pm
1710$\,ng/dL. The gap is cohort-level, not error: Ichihara's own
$\mathrm{CL}_{\text{tot}}/F$ ($14.3$\,L/day/kg $\approx 930$\,L/day at
65\,kg) is about half the $\mathrm{CL}_T = 1817.4$\,L/day inherited from
Bi, and their $V_d/F$ ($118.7$\,L/kg $\approx 7700$\,L) about half Bi's
14\,400\,L. Adopting Ichihara's clearance for TE would put TE and TC
$2\times$ apart on the \emph{same molecule}, breaking the
conservation-of-mass identity. Magnitude is therefore anchored to Bi;
only shape was taken from Ichihara. The per-user calibration factor
$\alpha$ exists to absorb individual-level differences of this
magnitude, and Ichihara's $C_{\max}$ CV of 0.58 confirms that
between-subject spread is wide enough to accommodate the gap.

\subsubsection{Testosterone undecanoate (TU)}

TU's terminal half-life (20.0\,d) is retained from Zhang et
al.~\cite{zhang1998}, who measured $18.3 \pm 2.3$\,d (500\,mg) and
$23.7 \pm 2.7$\,d (1000\,mg); 20.0 sits inside both arms. The absorption
rate is fitted so that $t_{\max}$ lands on the Aveed (FDA) clinical
pharmacology median of 7\,d (range 4--42\,d): solving
Equation~\ref{eq:tmax} with $k_a = \ln 2 / 20.0 = 0.0347$\,day$^{-1}$
for $k_e$ gives $k_e = 0.3748$\,day$^{-1}$ ($t_{1/2} = 44.4$\,h), and
$t_{\max} = 7.00$\,d.

The scaling constant follows from the identity:
$S = 10^5 \times 0.3748 / 1817.4 = 20.62$.

\paragraph{Peak/trough ratio.} Over a 70-day interval, the model
produces a peak/trough ratio of 8.1$\times$ (1045/129\,ng/dL). Wang et
al.~\cite{wang2010} report a steady-state peak/trough of 1.8$\times$
(789.8/435.6\,ng/dL at week 12). This gap is \emph{arithmetic, not
structural}: over a 70-day interval, a terminal half-life of 20\,d is
3.5 half-lives of decay, so peak/trough cannot fall below $\sim$11$\times$
on decay alone. Both targets \emph{are} jointly reachable in the
two-rate Bateman---a 73.8\,d terminal half-life paired with a 27.5\,h
release rate hits $t_{\max} = 7$\,d \emph{and} peak/trough 1.80$\times$
simultaneously---but 73.8\,d is 3.7$\times$ what Zhang measured. The
blocker is that Zhang's measured terminal half-life and Wang's observed
flatness cannot both be true; resolving it requires Wang's primary
sampling schedule, not a model rewrite. This finding is consistent with
Pastuszak et al.~\cite{pastuszak2021}, who fit injectable TU with a
one-compartment first-order-absorption model---the same two-rate-constant
Bateman used here.

\subsubsection{Testosterone propionate (TP)}

TP remains uncalibrated. It retains a placeholder 2\,h $k_e$
(physiological clearance of free testosterone), which pins its
$t_{\max}$ to 0.33\,d against a published window of 1--2\,d. Its
\emph{magnitude} was corrected by the clearance identity
(Section~\ref{sec:identity}), so steady-state $C_{\mathrm{avg}}$ is
correct, but the first $\sim$24\,h of the curve---where $t_{\max}$
lives---is wrong. The best available source is a textbook
chapter~\cite{behre2004} showing a simulated twice-weekly curve rather
than fitted rate constants or raw single-dose data. A fabricated fit
would be worse than the honest flag.

\subsection{Estradiol ester derivations}
\label{sec:e2}

\subsubsection{Estradiol valerate (EV) and cypionate (EC)}

EV and EC are calibrated against transfemscience.org's compiled
tables~\cite{tfs}, which aggregate published single-dose and
steady-state estradiol PK. The calibration procedure:
\begin{enumerate}[nosep,leftmargin=1.4em]
\item $k_e$ was fit so Equation~\ref{eq:tmax} reproduces the published
  time-to-peak (2.1\,d for EV, 3.1--4.3\,d for EC).
\item $k_a$ was fit so the modeled steady-state trough reproduces
  Table 10's 5\,mg-weekly figures.
\item $S$ was anchored to the same table's steady-state peak.
\end{enumerate}

The resulting constants reproduce the published steady-state $C_{\max}$
and $C_{\min}$ to within 0.1\% (Table~\ref{tab:accuracy}). A key
decision: single-dose and steady-state targets were \emph{not
simultaneously satisfiable} in the two-parameter model, and steady state
was chosen on the reasoning that it is what users actually live at and
what clinicians titrate on. The consequence is that single-dose
$C_{\max}$ is a derived value, not an enforced target
(Section~\ref{sec:anchor_choice}).

\subsubsection{Oral estradiol}

Oral estradiol uses the same Bateman biexponential with fast rate
constants. The absorption rate $k_a = 8.762$\,day$^{-1}$
($t_{1/2} = 0.0791$\,d) was fit so that $t_{\max} = 6.0$\,h, the
midpoint of the published 4--8\,h window. The elimination rate $k_e =
1.386$\,day$^{-1}$ ($t_{1/2} = 12.0$\,h) and $S = 1250$ were anchored to
reproduce the steady-state average for 2\,mg daily (90\,pg/mL).

A coupling issue arises: $S$ and $k_a$ are coupled through the peak but
not through the average (integrating the Bateman term gives
$\mathrm{AUC}/\text{dose} = D^{\mathrm{eff}} S / k_e$, with no $k_a$ in
it). One $(S, k_a)$ pair cannot hold both the published $t_{\max}$ and
the published single-dose $C_{\max}$. Anchoring on $C_{\mathrm{avg}}$
keeps 2\,mg daily at 90\,pg/mL---unchanged by the $k_a$ refit, and the
quantity clinicians titrate on---at the cost of single-dose $C_{\max}$
reading 88\,pg/mL rather than 60\,pg/mL for 2\,mg.

\subsubsection{Transdermal estradiol}

Transdermal estradiol uses the 24 hourly micro-dose construction
(Equation~\ref{eq:transdermal}) with $S = 155$, $F = 0.15$, and $k_a =
\ln 2 / 1.5 = 0.462$\,day$^{-1}$. The stored $k_e$ is unused because the
micro-dose branch sets $k_e = 0$ (single-exponential decay).

\subsection{Transdermal testosterone}
\label{sec:td_t}

Anchored to the AndroGel 1\% prescribing information~\cite{androgel},
Table 5 (Day 180 steady state, $n=44$): 50\,mg/day $\to
C_{\mathrm{avg}} = 555$\,ng/dL. The terminal half-life comes from the
label's washout statement---levels return to baseline in 48--72\,h,
implying $\sim$5\% residual at 60\,h, giving $t_{1/2} = 0.5784$\,d.
Transdermal $F = 0.10$ is stated outright in the label's \S12.3:
``Approximately 10\% of the testosterone dose applied on the skin
surface from AndroGel is absorbed into systemic circulation.''

The scaling constant $S = 132.98$ is derived from the 50\,mg anchor.
Two limitations are documented rather than fitted around; both are in
Section~\ref{sec:limitations}.

\subsection{Oral testosterone undecanoate}
\label{sec:oral_t}

Anchored to Swerdloff et al.~\cite{swerdloff2020} / the Jatenzo label:
237\,mg twice daily $\to C_{\mathrm{avg}} = 488.7$\,ng/dL, $t_{\max}$
2--4\,h. Both rate constants were solved jointly so that $t_{\max}$
lands at 3\,h and the steady-state $C_{\max}/C_{\mathrm{avg}}$ comes as
close as the model permits to the reported 2.06. The resulting constants
are $k_a = 7.702$\,day$^{-1}$ ($t_{1/2} = 0.09$\,d) and $k_e =
8.305$\,day$^{-1}$ ($t_{1/2} = 2.003$\,h), with $S = 13.57$.

\subsection{Dynamic endogenous baseline}
\label{sec:baseline}

Endogenous hormone production is modeled as a time-varying baseline
$C_{\text{base}}(t)$ that decays from an intact gonadal level to a
suppressed level as HRT progresses. The model uses a saturating
exponential taper:
\begin{equation}
C_{\text{base}}(t) = C_{\text{full}} + (C_{\text{supp}} - C_{\text{full}}) \cdot \left(1 - e^{-2.5 \, p(t)}\right)
\label{eq:baseline}
\end{equation}
where $p(t) = \min(\text{weeks on HRT} / t_{\text{scale}}, 1)$ is the
normalized progress, $t_{\text{scale}}$ is user-configurable (4, 8, or
12 weeks), $C_{\text{full}}$ is the intact baseline (100\,pg/mL for
estradiol, 600\,ng/dL for testosterone), and $C_{\text{supp}}$ is the
suppressed baseline (5\,pg/mL for estradiol, 15\,ng/dL for
testosterone). The factor 2.5 in the exponent ensures the suppression
reaches $\sim$92\% at $p = 1$. After 104 weeks (2 years), the baseline
is treated as fully suppressed.

\paragraph{Anti-androgen override.} When the user reports taking an
anti-androgen (spironolactone, cyproterone, or a GnRH agonist), the
testosterone baseline is set to the suppressed level immediately,
reflecting the direct gonadal suppression these agents provide.

\subsection{Blood test calibration}
\label{sec:calibration}

\subsubsection{The calibration factor}

Each blood test record stores a measured hormone level and the date the
blood was drawn. The calibration factor is the ratio of measured to
predicted concentration at the draw instant:
\begin{equation}
\alpha = \frac{C_{\text{measured}}}{C_{\text{predicted}}}
\label{eq:alpha}
\end{equation}
where $C_{\text{predicted}}$ is computed by running the PK model at the
exact draw date \emph{without} calibration applied (to avoid
circularity: the calibration engine must never feed on its own output).
The factor is clamped to $[0.2, 5.0]$ to exclude clinically implausible
values. When multiple calibrated blood tests exist, the most recent one
for the currently selected hormone type is used.

\subsubsection{Application}

$\alpha$ is applied as a multiplicative scalar to the entire model
output (Equation~\ref{eq:superposition}). This is a pragmatic per-user
analog of the inter-individual variability (IIV) terms in a formal
mixed-effects population-PK model. The published literature reports
substantial between-subject variability: Ichihara et al.~\cite{ichihara2020}
report a CV of 0.58 on $V_d/F$ for TE, corresponding to an $\sim$8-fold
individual range. The clamp range $[0.2, 5.0]$ spans most of this
distribution without admitting values that would indicate a data entry
error.

\subsubsection{Hormone-type gating}

The calibration factor is only applied when the blood test's hormone
type matches the currently selected medication. This prevents a
testosterone draw from silently rescaling the estradiol curve, which
would occur if the filter were absent.

\subsection{Drug interactions: spironolactone}
\label{sec:spiro}

Spironolactone is a known CYP3A4 inducer that can modestly increase
estradiol clearance. When the user has opted in to the spironolactone
clearance adjustment, the elimination rate $k_e$ for estrogen doses is
scaled by a factor of 1.15 (15\% faster clearance):
\begin{equation}
k_e' = k_e \cdot 1.15
\label{eq:spiro}
\end{equation}
This adjustment is applied at the PK parameter resolution level, before
the Bateman function is evaluated. It is not applied to testosterone
doses, as the interaction is specific to estradiol metabolism.

Note that because $t_{\max}$ depends on $k_e$
(Equation~\ref{eq:tmax}), this adjustment also shifts the peak time
slightly earlier. The ester-level $t_{\max}$ values reported in
Table~\ref{tab:tmax} are computed without the spironolactone adjustment.

\subsection{Aromatization estimation}
\label{sec:aromatization}

For trans men on testosterone therapy, a small fraction of circulating
testosterone is converted to estradiol via aromatase. The model
estimates this conversion as:
\begin{equation}
E_2^{\text{arom}} = 0.002 \times T_{\text{ng/dL}} \quad \text{(pg/mL)}
\label{eq:aromatization}
\end{equation}
i.e., a 0.2\% conversion rate. This helps trans men monitor whether
their estradiol is within the male reference range (10--40\,pg/mL). The
conversion factor is a population average; individual aromatase activity
varies with body composition, age, and genetic factors.

\subsection{Chart window and sampling resolution}
\label{sec:chart_window}

The chart's visible time range and sampling step are derived from the
ester's pharmacokinetic properties, not hardcoded. The rule is:
\begin{align}
W    &= \mathrm{clamp}(1.5 \cdot t_{\max},\; \text{lower:}\;\min(2\,\text{d},\; 2\tau),\; \text{upper:}\;1.5\tau) \label{eq:window} \\
\text{step} &= \max\bigl(\min(W/40,\; t_{\max}/8),\; 15\,\text{min}\bigr) \label{eq:step}
\end{align}
where $W$ is the total window width, $\tau$ is the dosing interval, and
the window is split 1:2 past:future. The $t_{\max}/8$ cap is an
anti-aliasing guarantee: it ensures the sampling step never exceeds
one-eighth of the time-to-peak, preventing the drawn curve from
collapsing into a flat line when the step would otherwise land on the
same phase of every dose cycle. The $W/40$ term keeps every ester at a
comparable point count ($\sim$41) so curves read with the same density
regardless of timescale.

When $t_{\max}$ is uncalibrated or non-finite (the transdermal routes,
which have no meaningful $t_{\max}$ because $k_e = 0$), a fallback
window of 2\,d is used.

% ======================================================================
\section{Encoded constants}
\label{sec:constants}

\begin{table}[h]
\centering
\small
\begin{tabular}{@{}lrrrrrrr@{}}
\toprule
& \multicolumn{2}{c}{\textbf{slow exponential}} & \textbf{fast} & & & & \\
\cmidrule(lr){2-3}
\textbf{Form} & $t_{1/2}$ (d) & $k_a$ (d$^{-1}$) & $k_e$ (d$^{-1}$)
               & EWF & $F$ & $S$ & $t_{\max}$ \\
\midrule
EV            & 2.6000 & 0.2667 & 0.7737 & 0.764 & 1.00 & 396.40 & 2.10\,d \\
EC            & 6.4700 & 0.1071 & 0.4310 & 0.687 & 1.00 & 273.00 & 4.30\,d \\
E2 oral       & 0.0791 & 8.7629 & 1.3859 & 1.000 & 0.05 & 1250.00 & 6.0\,h \\
E2 cream      & 1.5000 & 0.4621 & ---    & 1.000 & 0.15 & 155.00 & --- \\
\midrule
TC            & 3.8390 & 0.1806 & 1.2196 & 0.699 & 1.00 & 67.13 & 1.84\,d \\
TE            & 6.6000 & 0.1050 & 1.9625 & 0.720 & 1.00 & 107.98 & 1.58\,d \\
TP            & 1.0000 & 0.6931 & 8.3178 & 0.837 & 1.00 & 457.67 & 0.33\,d \\
TU            & 20.0000 & 0.0347 & 0.3748 & 0.631 & 1.00 & 20.62 & 7.00\,d \\
T oral (TU)   & 0.0900 & 7.7016 & 8.3053 & 0.631 & 1.00 & 13.57 & 3.0\,h \\
T gel         & 0.5784 & 1.1984 & ---    & 1.000 & 0.10 & 132.98 & --- \\
\bottomrule
\end{tabular}
\caption{Constants as shipped, and the $t_{\max}$ each implies via
Equation~\ref{eq:tmax}. The cream/gel routes have no $t_{\max}$ because
they are not modeled with Equation~\ref{eq:bateman}; their $k_e$ is
stored but unused.}
\label{tab:constants}
\end{table}

\paragraph{Calibration status.} Of the eight forms with a $t_{\max}$,
seven are calibrated against a primary or label source: EV and EC
against compiled steady-state tables~\cite{tfs}, oral E2 against the
published 4--8\,h window, TC against Bi et al.~\cite{bi2018}, TE against
Ichihara et al.~\cite{ichihara2020}, TU against Aveed's FDA clinical
pharmacology, and oral T against the Jatenzo label~\cite{swerdloff2020}.
TP alone remains uncalibrated: its $t_{\max}$ of 0.33\,d is known to be
wrong against a published window of 1--2\,d, and is not presented to
users as a fact about the drug.

% ======================================================================
\section{Results}
\label{sec:results}

\subsection{Accuracy against anchors}
\label{sec:accuracy}

\begin{table}[h]
\centering
\small
\begin{tabular}{@{}lllrrl@{}}
\toprule
\textbf{Regimen} & \textbf{Quantity} & \textbf{Model} & \textbf{Published}
                 & \textbf{Error} & \textbf{Source} \\
\midrule
EV 5\,mg q7d & $C_{\max}$ & 384 & 384 & \good{$+0.1\%$} & TFS Table 10 \\
EV 5\,mg q7d & $C_{\min}$ & 142 & 142 & \good{$+0.1\%$} & TFS Table 10 \\
EC 5\,mg q7d & $C_{\max}$ & 339 & 339 & \good{$-0.1\%$} & TFS Table 10 \\
EC 5\,mg q7d & $C_{\min}$ & 262 & 262 & \good{$-0.1\%$} & TFS Table 10 \\
TC 100\,mg q7d & $C_{\mathrm{avg}}$ & 550 & 549 & \good{$+0.0\%$} & Bi Table 2 \\
T gel 50\,mg/d & $C_{\mathrm{avg}}$ & 555 & 555 & \good{$+0.1\%$} & AndroGel Table 5 \\
T gel 75\,mg/d & $C_{\mathrm{avg}}$ & 833 & 601 & \bad{$+38.6\%$} & AndroGel Table 5 \\
T gel 100\,mg/d & $C_{\mathrm{avg}}$ & 1111 & 713 & \bad{$+55.8\%$} & AndroGel Table 5 \\
Oral TU 237\,mg BID & $C_{\mathrm{avg}}$ & 489 & 489 & \good{$-0.0\%$} & Jatenzo \\
Oral TU 237\,mg BID & $C_{\max}$ & 789 & 1008 & \bad{$-21.7\%$} & Jatenzo \\
\bottomrule
\end{tabular}
\caption{Predicted versus published. Rows with near-zero error are the
anchors themselves and demonstrate that the fit converged. The three
coloured rows are the informative ones; both are explained in
Section~\ref{sec:limitations}. Estradiol in pg/mL, testosterone in
ng/dL.}
\label{tab:accuracy}
\end{table}

Rows in Table~\ref{tab:accuracy} with 0.0\% error demonstrate only that
the fit converged---a fit agrees with its own target by construction.
The three coloured rows are the informative ones.

\subsection{Out-of-sample validation}
\label{sec:out_of_sample}

The following comparisons were \emph{not} used in any fit.

\begin{table}[h]
\centering
\footnotesize
\begin{tabular}{@{}lrrrl@{}}
\toprule
\textbf{Check} & \textbf{Model} & \textbf{Published} & \textbf{Error} & \textbf{Source} \\
\midrule
TC single 200\,mg $C_{\max}$ (exogenous) & 997 & 872 & $+14.3\%$
  & 200\,mg IM lit.\ 1112 less 240 basal \\
TC 100\,mg q7d $C_{\max}$ + basal & 976 & 1150 & $-15.1\%$
  & Bi Figure 2, 100\,mg panel \\
EV single 2\,mg $C_{\max}$ & 119 & 118 & \good{$+1.0\%$} & TFS Table 9 \\
EC single 2\,mg $C_{\max}$ & 59 & 62 & \good{$-5.1\%$} & TFS Table 9 \\
\bottomrule
\end{tabular}
\caption{Independent validation. The estradiol rows are the strongest
evidence: constants fit purely to steady-state data reproduce
single-dose peaks within 1\% and 5\%.}
\label{tab:out_of_sample}
\end{table}

The two testosterone rows agree to about 15\% in opposite directions,
which is roughly the precision the source itself supports given Bi's
reported RSEs. They are consistency checks, not tight validation.

The estradiol rows are the strongest evidence in the model: constants
fit purely to steady-state data reproduce single-dose peaks within 1\%
and 5\%. This is a non-trivial result---it demonstrates that the
two-parameter Bateman, with steady-state-anchored constants, generalizes
to the single-dose regime without refitting.

\subsection{Time to peak}
\label{sec:tmax_results}

\begin{table}[h]
\centering
\small
\begin{tabular}{@{}llll@{}}
\toprule
\textbf{Form} & \textbf{Model $t_{\max}$} & \textbf{Published} & \textbf{Source quality} \\
\midrule
EV      & 2.10\,d & 2.1\,d & compiled tables \\
EC      & 4.30\,d & 3.1--4.3\,d & compiled tables \\
E2 oral & \good{6.0\,h} & \good{4--8\,h} & secondary \\
TC      & \good{1.84\,d} & \good{1.84\,d} & \textbf{primary} (Bi Table 2) \\
TE      & \good{1.58\,d} & \good{1.7 $\pm$ 0.9\,d} & \textbf{primary} (Ichihara 2020) \\
TP      & \bad{0.33\,d} & \bad{1--2\,d} & sparse (uncalibrated) \\
TU      & \good{7.00\,d} & \good{median 7\,d} & \textbf{label} (Aveed FDA) \\
T oral  & \good{3.0\,h} & \good{2--4\,h} & label (Jatenzo) \\
\bottomrule
\end{tabular}
\caption{Where the model puts the peak. Seven of eight forms with a
$t_{\max}$ are calibrated against a primary or label source. TP alone
remains uncalibrated: its $t_{\max}$ is known to be wrong and is not
presented to users as fact.}
\label{tab:tmax}
\end{table}

\subsection{Persona validation}
\label{sec:personas}

Section-level agreement with published tables is necessary but not
sufficient: those tables describe clean steady state on a fixed regimen,
and real users miss doses, change dose, change route, and dose on
intervals nobody studied. Eight synthetic users were run end to end
through the same summation the app performs.

\begin{longtable}{@{}p{2.6cm}p{4.5cm}p{7.4cm}@{}}
\toprule
\textbf{Persona} & \textbf{Regimen} & \textbf{What it tests, and the result} \\
\midrule
\endhead
Maya, 24 & 4\,mg EV IM q7d, 18 months, perfect adherence
 & Control. $C_{\max}$ 307, $C_{\min}$ 114\,pg/mL --- exactly $4/5$ of
   the 5\,mg source figures, as linearity requires. \good{Pass.} \\[0.4em]
Devon, 31 & 100\,mg TC q7d $\to$ 200\,mg q14d at month 6
 & Superposition across a regimen change. $C_{\mathrm{avg}}$ preserved
   at 790\,ng/dL as it must be; trough falls $561 \to 381$ while peak
   rises $976 \to 1339$, and peak/trough widens $1.74\times \to
   3.51\times$. The model reproduces the clinical trade-off and
   correctly refuses to reward the longer interval. \good{Pass.} \\[0.4em]
Priya, 47 & 2\,mg oral E2 BID for 4 months, then 4\,mg EV q7d
 & Route switch; historical doses must keep their own kinetics. They
   do --- the oral segment retains $F = 0.05$ and the fast rates.
   Modeled BID swing is $1.37\times$ at an unchanged $C_{\mathrm{avg}}$
   of 180\,pg/mL. \good{Pass.} \\[0.4em]
Sam, 19 & 50\,mg/day T gel, misses $\sim$2 days/week
 & Skipped-dose handling. $C_{\mathrm{avg}}$ 637\,ng/dL; missed days
   cost 20\% of average and 63\% of trough. \good{Pass.} \\[0.4em]
Rowan, 38 & 4\,mg/day compounded E2 cream, onboarded in-app
 & Route pseudo-ester resolution. Cream dose correctly resolves to the
   transdermal pseudo-ester constants. \good{Pass.} \\[0.4em]
Kai, 52 & 237\,mg oral TU BID, 10-day interruption
 & Washout and re-accumulation on the fastest route. Falls to the
   endogenous floor within a day; no phantom drug survives the gap;
   recovers to 100\% of pre-gap within 2 days. \good{Pass.} \\[0.4em]
Jules, 28 & 100\,mg TE IM q7d
 & TE steady-state. $C_{\max}$ 703, $C_{\min}$ 405, $C_{\mathrm{avg}}$
   566\,ng/dL (exogenous). Consistent with TC via the clearance
   identity. $t_{\max}$ 1.58\,d, inside Ichihara's observed $1.7 \pm
   0.9$\,d. \good{Pass.} \\[0.4em]
Nadia, 35 & 5\,mg EC q10d --- an interval nobody studied
 & Extrapolation outside the calibration envelope. Trough holds in
   band but sits 40\% below the same dose at q7d. Structurally sound
   but inherits every assumption of the fit. \good{Pass, with caveat.} \\
\bottomrule
\end{longtable}

All eight land inside their clinical reference bands. The persona suite
tests properties that section-level accuracy tables cannot: linear
scaling (Maya), superposition across regimen changes (Devon), route
switching with per-dose kinetics (Priya), skipped-dose handling (Sam),
route pseudo-ester resolution (Rowan), washout and re-accumulation
(Kai), inter-ester consistency (Jules), and extrapolation outside the
calibration envelope (Nadia).

% ======================================================================
\section{Discussion}
\label{sec:discussion}

\subsection{Structural limitations}
\label{sec:limitations}

These are quantified rather than fixed. Each needs either a source the
model does not have or a structurally different approach.

\subsubsection{Testosterone undecanoate peak/trough ratio}
\label{sec:tu_limitation}

Over a 70-day interval, the model produces a peak/trough ratio of
8.1$\times$ against Wang et al.'s~\cite{wang2010} observed 1.8$\times$.
This is not a structural deficiency of the two-rate Bateman: as noted in
Section~\ref{sec:te_tu}, both targets are jointly reachable with a
73.8\,d terminal half-life, but that value is 3.7$\times$ what Zhang et
al.~\cite{zhang1998} measured. The blocker is a \emph{source conflict}
between Zhang's terminal half-life and Wang's observed flatness, not the
number of compartments. Pastuszak et al.~\cite{pastuszak2021} fit
injectable TU with the same one-compartment first-order-absorption model
used here, confirming the model family is appropriate.

The magnitude correction via the clearance identity improves TU's
$C_{\mathrm{avg}}$ (496\,ng/dL exogenous, against a published
$\sim$500--600 total) but its $C_{\max}$ overshoots (1045 vs.\ 790\,ng/dL
observed), because correcting magnitude cannot correct shape. The trade
was taken knowingly: $C_{\mathrm{avg}}$ is what clinicians titrate on
and what a user lives at.

\subsubsection{Transdermal intraday ripple}
\label{sec:td_ripple}

The 24-hourly-micro-dose construction has no absorption phase and tops
out at $C_{\max}/C_{\min} \approx 1.05$ for testosterone gel, against an
observed 2.24. Reproducing 2.24 would require a terminal half-life near
1\,h, which contradicts the label's own 48--72\,h washout by an order of
magnitude. The two cannot be satisfied together in this model.
$C_{\mathrm{avg}}$ is exact.

\subsubsection{Transdermal dose-response sublinearity}
\label{sec:td_dose_response}

The AndroGel 1\% label's dose-response is markedly sublinear:
50\,mg $\to$ 555, 75\,mg $\to$ 601, 100\,mg $\to$ 713\,ng/dL. The 1.62\%
label inverts outright: 81\,mg $C_{\mathrm{avg}}$ 537 is \emph{below}
60.75\,mg $C_{\mathrm{avg}}$ 643.

This is a titration artifact: higher-dose arms are enriched for poor
absorbers, because subjects were escalated precisely when they responded
badly. It is selection confounding, not saturable pharmacokinetics, and
a linear model should not be bent to reproduce it. $S$ is anchored to
the 50\,mg arm---the standard starting dose and the least
confounded---so predictions at 75 and 100\,mg run high by 39\% and 56\%
(Table~\ref{tab:accuracy}).

Gel absorption also varies enormously between individuals. The
calibration factor $\alpha$ is the intended remedy and is more important
on this route than on any other.

\subsubsection{Oral testosterone peak fluctuation}
\label{sec:oral_t_limitation}

With $t_{\max}$ pinned to the reported 2--4\,h window, the model's
achievable $C_{\max}/C_{\mathrm{avg}}$ tops out at 1.61 against a
reported 2.06, so $C_{\max}$ reads 789 rather than 1008\,ng/dL.

The reported ratio is itself inflated by mean-of-maxima bias: the
published $C_{\max}$ is the mean of per-subject maxima, and
$\mathbb{E}[\max] \ge \max[\mathbb{E}]$, with a reported SD of $\pm 581$
on a mean of 1008. Anchoring to $C_{\mathrm{avg}}$---the label's own
primary endpoint---and accepting a documented under-prediction of the
peak was judged the better trade.

\subsubsection{Testosterone propionate remains uncalibrated}
\label{sec:tp_limitation}

TP retains a placeholder 2\,h $k_e$, which pins its $t_{\max}$ to 0.33\,d
against a published window of 1--2\,d. Its magnitude was corrected by
the clearance identity, so steady-state $C_{\mathrm{avg}}$ is correct,
but the peak arrives too early and the curvature between peak and trough
is wrong. The best available source is a textbook chapter~\cite{behre2004}
showing a simulated curve, not fitted rate constants. A fabricated fit
would be worse than the honest flag.

\subsubsection{Model-wide simplifications}
\label{sec:simplifications}

The model assumes:
\begin{itemize}[nosep,leftmargin=1.4em]
\item \textbf{Dose-linearity} with no saturable protein binding or SHBG
  dynamics.
\item \textbf{A single well-mixed compartment} with constant clearance.
\item \textbf{No covariates}: body composition, age, hepatic function,
  and concomitant medications (beyond the one spironolactone adjustment)
  are not represented.
\item \textbf{No between-subject variability} in the population
  prediction. Predictions are population means, not individual
  forecasts. The per-user calibration factor $\alpha$ exists to close
  some of that gap once blood tests are available.
\end{itemize}

These are standard simplifications for a one-compartment model and are
shared by the published population-PK models cited
above~\cite{bi2018,ichihara2020,pastuszak2021}, with the caveat that
those models additionally include IIV terms and covariates (e.g., body
weight and SHBG in Pastuszak's TU model).

\subsection{The steady-state vs.\ single-dose anchor choice}
\label{sec:anchor_choice}

For both estradiol esters and all four testosterone esters, the
single-dose peak targets and the published steady-state data were
\emph{mutually unsatisfiable} in the two-parameter model. For estradiol
the gap was large: holding the single-dose targets put steady state
$+110\%$ (EV) and $+206\%$ (EC) above the published table.

Steady state was chosen, in every case, on the reasoning that it is what
users actually live at and what clinicians titrate on. The consequence
is that single-dose figures are derived values, not enforced targets.
This is a modelling decision, not a derivation, and it has a cost:
single-dose $C_{\max}$ for EV (119\,pg/mL for 2\,mg) and EC (59\,pg/mL
for 2\,mg) are lower than some published single-dose figures, though
they agree with the transfemscience Table 9 values to within 1\% and
5\% respectively (Table~\ref{tab:out_of_sample}).

\subsection{Source quality hierarchy}
\label{sec:source_quality}

The model's constants rest on a hierarchy of source quality:

\begin{enumerate}[nosep,leftmargin=1.4em]
\item \textbf{Primary population-PK} (TC): Bi et al.~\cite{bi2018}
  supply $\mathrm{CL}/F$, $V/F$, and $K_a$ with RSEs and bootstrap CIs.
  This is the strongest source in the model.
\item \textbf{Primary single-dose PK} (TE): Ichihara et
  al.~\cite{ichihara2020} supply terminal half-life and per-subject
  absorption rates in the target population (transgender men, IM
  route).
\item \textbf{Regulatory labels} (T gel, oral T, TU $t_{\max}$):
  AndroGel~\cite{androgel} and Jatenzo~\cite{swerdloff2020} supply
  steady-state $C_{\mathrm{avg}}$, $C_{\max}$, and $t_{\max}$ from FDA
  registration trials. Aveed's FDA clinical pharmacology supplies TU's
  median $t_{\max}$.
\item \textbf{High-quality secondary compilation} (EV, EC):
  transfemscience.org~\cite{tfs} aggregates published single-dose and
  steady-state estradiol data. Not a primary table, and flagged as
  such.
\item \textbf{Secondary literature} (TU terminal half-life): Zhang et
  al.~\cite{zhang1998} and Wang et al.~\cite{wang2010} supply
  terminal half-life and steady-state peak/trough, respectively. These
  two sources are in conflict (Section~\ref{sec:tu_limitation}).
\item \textbf{Nothing} (TP): beyond the clearance identity that ties
  TP's magnitude to TC, there is no primary source for TP's absorption
  rate. The placeholder is flagged.
\end{enumerate}

\subsection{The per-user calibration factor as a pragmatic IIV analog}
\label{sec:iiv}

The published population-PK models for testosterone esters all report
substantial inter-individual variability. Ichihara et al.~\cite{ichihara2020}
report a CV of 0.58 on $V_d/F$ for TE, and Bi et al.~\cite{bi2018} report
IIV terms on clearance and volume. A formal mixed-effects model would
represent this variability with $\eta$ terms on each parameter, drawn
from a population distribution.

xT's calibration factor $\alpha$ is a simpler approach: a single scalar
multiplier on the entire model output, derived from the user's own blood
test. It cannot represent parameter-specific variability (e.g., a user
with fast clearance but normal volume), and it cannot provide
population-level confidence intervals. But it is practical for a
consumer application: it requires no raw data access, no fitting
infrastructure, and no statistical expertise from the user, and it
absorbs the dominant source of individual-level error once a blood test
is available.

The clamp range $[0.2, 5.0]$ spans most of the reported IIV distribution
without admitting values that would indicate a data entry error. For
routes with especially high between-subject variability---transdermal in
particular---$\alpha$ is more important than on any other route.

% ======================================================================
\section{Conclusion}
\label{sec:conclusion}

We have presented the pharmacokinetic model underlying xT, a mobile
application for hormone replacement therapy monitoring. The model
employs a one-compartment Bateman biexponential with linear
superposition of dose contributions, a dynamic endogenous baseline, and
a per-user calibration factor. Ten medication/route/form combinations
are supported across estradiol and testosterone.

The model's key structural features are:
\begin{itemize}[nosep,leftmargin=1.4em]
\item A \textbf{conservation-of-mass identity} that constrains the
  scaling constant $S$ for all injectable testosterone esters to depend
  only on $k_e$ and clearance, eliminating $S$ as a free parameter and
  enforcing mutual consistency across esters.
\item An \textbf{analytic closed-form derivation} of $S$ for
  testosterone cypionate from primary population-PK parameters, which
  anchors the identity.
\item A \textbf{per-user calibration factor} $\alpha$ that serves as a
  pragmatic analog of inter-individual variability terms in formal
  mixed-effects models.
\item A \textbf{dynamic endogenous baseline} that models gonadal
  suppression over time, separating exogenous from endogenous
  contributions.
\end{itemize}

The model achieves 0.0--0.1\% accuracy against its anchors and 1--15\%
on out-of-sample checks. Seven of eight forms with a $t_{\max}$ are
calibrated against primary or label sources; one (TP) is honestly
flagged as uncalibrated. Structural limitations are quantified rather
than hidden: TU's peak/trough ratio is a source conflict, not a
structural deficiency; transdermal intraday ripple is under-modeled by
construction; and oral peak fluctuation is under-predicted by an
anchoring choice that prioritizes $C_{\mathrm{avg}}$.

The model is honest about where its constants are primary-sourced, where
they are secondary, and where they remain uncalibrated. This honesty---
the explicit marking of what is known, what is estimated, and what is
placeholder---is, in our view, the most important property of a
pharmacokinetic model used in a consumer health context, where users
may make medical decisions based on its output.

\vfill
\noindent\rule{\linewidth}{0.4pt}\\[0.3em]
{\footnotesize This paper describes a consumer tracking application's
internal model. It is not clinical guidance, and its predictions are
population estimates that should not be used to make dosing decisions
without laboratory confirmation.}

% ======================================================================
\begin{thebibliography}{99}

\bibitem{bi2018}
Bi Y., et al.\ (2018). ``Population Pharmacokinetic/Pharmacodynamic
Modeling of Depot Testosterone Cypionate in Healthy Male Subjects.''
\emph{CPT: Pharmacometrics \& Systems Pharmacology}.
\texttt{doi:10.1002/psp4.12287}.

\bibitem{ichihara2020}
Ichihara K., Masumori N., Fujii S., Toda T.\ (2020). ``Pharmacokinetics
of Testosterone Enanthate After Intramuscular Injection for Transgender
Men.'' \emph{Androgens: Clinical Research and Therapeutics}.
\texttt{doi:10.1089/andro.2020.0002}.

\bibitem{zhang1998}
Zhang G.Y., Gu Y.Q., Wang X.H., Cui Y.G., Bremner W.J.\ (1998). ``A
pharmacokinetic study of injectable testosterone undecanoate in
hypogonadal men.'' \emph{J.\ Androl.} 19(6):761--8.
PMID: 9876028.

\bibitem{wang2010}
Wang C., et al.\ (2010). ``Pharmacokinetics and Safety of Long-Acting
Testosterone Undecanoate Injections in Hypogonadal Men: An 84-Week
Phase III Clinical Trial.'' \emph{J.\ Androl.}
\texttt{doi:10.2164/jandrol.109.009597}.

\bibitem{pastuszak2021}
Pastuszak A.W., Gittelman M., Tursi J.P., Jaffe J.S., Miner M.M.\
(2021). ``Population Pharmacokinetic Modeling and Simulations to
Evaluate a Potential Dose Regimen of Testosterone Undecanoate in
Hypogonadal Males.'' \emph{J.\ Clin.\ Pharmacol.}
PMID: 34269421.

\bibitem{swerdloff2020}
Swerdloff R.S., et al.\ (2020) / Jatenzo (testosterone undecanoate)
prescribing information. FDA label.

\bibitem{androgel}
AndroGel 1\% (testosterone gel) prescribing information. FDA label,
\S12.3 and Table 5.

\bibitem{tfs}
transfemscience.org compiled tables (9 and 10). Estradiol ester
pharmacokinetics. Secondary compilation of published single-dose and
steady-state data.

\bibitem{behre2004}
Behre H.M., Wang C., Handelsman D.J., Nieschlag E.\ (2004).
``Pharmacology of testosterone preparations.'' In \emph{Testosterone:
Action, Deficiency, Substitution}, Cambridge University Press.

\bibitem{hembree2017}
Hembree W.C., et al.\ (2017). ``Endocrine Treatment of Gender-Dysphoric
/Gender-Incongruent Persons: An Endocrine Society Clinical Practice
Guideline.'' \emph{J.\ Clin.\ Endocrinol.\ Metab.} 102(11):3869--3903.

\bibitem{rowland2011}
Rowland M., Tozer T.N.\ (2011). \emph{Clinical Pharmacokinetics and
Pharmacodynamics: Concepts and Applications}. 4th ed. Lippincott
Williams \& Wilkins.

\end{thebibliography}

\end{document}